\section{嘉道时期的市场、河流与边境世界}
嘉道危机通过盐、漕、河、债务、矿业、走私、商路和边贸连接不同地区。本节将这些节点视为市场与环境共同塑造的制度问题：水灾、疾病、交通与价格不是政治背景，而是决定文书是否能转化为行动的条件。
\subsection{1790—1799：区域网络、制度媒介与环境条件}
这些事件的先后可以由账本定位，但每一项影响都须回到具体区域、材料与行动者，不能被压缩为抽象的国家趋势。
\hypertarget{civic:EQ-1796-JIAQING-ACCESSION}{}\paragraph{嘉庆元年（1796年）}嘉庆帝正式即位，白莲教战争与财政压力成为初政核心问题。本事件需要放在市场价格、交通路线、水文环境、边贸和地方债务的相互制约中理解，不把中央决策与区域经济割裂开来。 核心取证为《清仁宗实录》嘉庆元年相关条；宫中朱批奏折、内阁大库题本与《清史稿》相关志传。行政文书显示的是机构如何把道路、港口、人口、宗教或案件转换为可处理的类别；没有后续县档或物质证据，不能由分类推定实际生活。账本提示的研究线索为宫廷政治、制度运作与中央—地方信息链研究。事件所处条件为禅位后的双重权力格局限制新帝自主空间。后续变化应限于嘉庆朝逐步重整人事和军费管理。来源限制：以嘉庆元年（1796年）的同期文书为定位起点；官修记录、奏报或外交文本服务特定机构立场，具体日期、规模、地方执行与对手视角须以独立材料复核。
\hypertarget{civic:EQ-1797-WHITE-LOTUS-WAR}{}\paragraph{嘉庆二年（1797年）}白莲教战争持续扩展，清军依赖乡勇、团练与地方资源。本事件需要放在市场价格、交通路线、水文环境、边贸和地方债务的相互制约中理解，不把中央决策与区域经济割裂开来。 核心取证为《清仁宗实录》嘉庆二年相关条；军机处档、战事方略及地方奏报。编年条目可以为行动建立相对次序，却通常不保存劳动过程、物价变动、建筑材料或非精英的叙述。账本提示的研究线索为战争动员、军需、战场暴力与地方社会研究。事件所处条件为正规军难以适应分散山地战。后续变化应限于战争改变地方武装和财政结构，后果超出战区。来源限制：以嘉庆二年（1797年）的同期文书为定位起点；官修记录、奏报或外交文本服务特定机构立场，具体日期、规模、地方执行与对手视角须以独立材料复核。
\hypertarget{civic:EQ-1798-MILITIA-MOBILIZATION}{}\paragraph{嘉庆三年（1798年）}朝廷加强对乡勇、团练和地方绅士资源的动员以应对白莲教战争。本事件需要放在市场价格、交通路线、水文环境、边贸和地方债务的相互制约中理解，不把中央决策与区域经济割裂开来。 核心取证为《清仁宗实录》嘉庆三年相关条；地方志、督抚奏折及地方档案。编年条目可以为行动建立相对次序，却通常不保存劳动过程、物价变动、建筑材料或非精英的叙述。账本提示的研究线索为地方档案、迁徙、宗教、族群与社会动员研究。事件所处条件为常备军效率不足与地形分散相互作用。后续变化应限于地方武装的形成既服务镇压，也改变地方权力关系。来源限制：以嘉庆三年（1798年）的同期文书为定位起点；官修记录、奏报或外交文本服务特定机构立场，具体日期、规模、地方执行与对手视角须以独立材料复核。
\hypertarget{civic:EQ-1799-HESHEN-PURGE}{}\paragraph{嘉庆四年（1799年）}乾隆帝去世后，嘉庆帝逮治和珅并清理其政治网络。本事件需要放在市场价格、交通路线、水文环境、边贸和地方债务的相互制约中理解，不把中央决策与区域经济割裂开来。 核心取证为《清仁宗实录》嘉庆四年相关条；宫中朱批奏折、内阁大库题本与《清史稿》相关志传。行政文书显示的是机构如何把道路、港口、人口、宗教或案件转换为可处理的类别；没有后续县档或物质证据，不能由分类推定实际生活。账本提示的研究线索为宫廷政治、制度运作与中央—地方信息链研究。事件所处条件为太上皇去世使嘉庆得以直接行使君权。后续变化应限于清算重塑宫廷财政与用人，但不等于解决结构性财政问题。来源限制：以嘉庆四年（1799年）的同期文书为定位起点；官修记录、奏报或外交文本服务特定机构立场，具体日期、规模、地方执行与对手视角须以独立材料复核。
\subsection{1800—1809：区域网络、制度媒介与环境条件}
这些事件的先后可以由账本定位，但每一项影响都须回到具体区域、材料与行动者，不能被压缩为抽象的国家趋势。
\hypertarget{civic:EQ-1800-WHITE-LOTUS-FINANCE}{}\paragraph{嘉庆五年（1800年）}白莲教战争军费、捐输与地方筹饷加重国家财政压力。本事件需要放在市场价格、交通路线、水文环境、边贸和地方债务的相互制约中理解，不把中央决策与区域经济割裂开来。 核心取证为《清仁宗实录》嘉庆五年相关条；户部奏销、盐政、河工或海关档案。行政文书显示的是机构如何把道路、港口、人口、宗教或案件转换为可处理的类别；没有后续县档或物质证据，不能由分类推定实际生活。账本提示的研究线索为财政账册、市场网络、灾害环境与区域经济研究。事件所处条件为长期作战消耗常规财政与漕运资源。后续变化应限于报销数字、捐输与实际地方负担需分项核验。来源限制：以嘉庆五年（1800年）的同期文书为定位起点；官修记录、奏报或外交文本服务特定机构立场，具体日期、规模、地方执行与对手视角须以独立材料复核。
\hypertarget{civic:EQ-1801-WHITE-LOTUS-SUPPRESSION}{}\paragraph{嘉庆六年（1801年）}清军在若干战区逐步压缩白莲教起事者活动空间。本事件需要放在市场价格、交通路线、水文环境、边贸和地方债务的相互制约中理解，不把中央决策与区域经济割裂开来。 核心取证为《清仁宗实录》嘉庆六年相关条；军机处档、战事方略及地方奏报。编年条目可以为行动建立相对次序，却通常不保存劳动过程、物价变动、建筑材料或非精英的叙述。账本提示的研究线索为战争动员、军需、战场暴力与地方社会研究。事件所处条件为地方武装与分区清剿策略加强。后续变化应限于战事结束的地区差异不能由中央凯报统一概括。来源限制：以嘉庆六年（1801年）的同期文书为定位起点；官修记录、奏报或外交文本服务特定机构立场，具体日期、规模、地方执行与对手视角须以独立材料复核。
\hypertarget{civic:EQ-1802-WAR-REFUGEE-RELIEF}{}\paragraph{嘉庆七年（1802年）}战区流民、复垦和赈济成为地方善后事务。本事件需要放在市场价格、交通路线、水文环境、边贸和地方债务的相互制约中理解，不把中央决策与区域经济割裂开来。 核心取证为《清仁宗实录》嘉庆七年相关条；地方志、督抚奏折及地方档案。编年条目可以为行动建立相对次序，却通常不保存劳动过程、物价变动、建筑材料或非精英的叙述。账本提示的研究线索为地方档案、迁徙、宗教、族群与社会动员研究。事件所处条件为长期战乱破坏村落、田地和交通。后续变化应限于赈济命令不自动等于救助覆盖或人口恢复。来源限制：以嘉庆七年（1802年）的同期文书为定位起点；官修记录、奏报或外交文本服务特定机构立场，具体日期、规模、地方执行与对手视角须以独立材料复核。
\hypertarget{civic:EQ-1803-FORBIDDEN-CITY-ATTACK}{}\paragraph{嘉庆八年（1803年）}陈德等闯入紫禁城，暴露京师警卫与内廷管理问题。本事件需要放在市场价格、交通路线、水文环境、边贸和地方债务的相互制约中理解，不把中央决策与区域经济割裂开来。 核心取证为《清仁宗实录》嘉庆八年相关条；宫中朱批奏折、内阁大库题本与《清史稿》相关志传。行政文书显示的是机构如何把道路、港口、人口、宗教或案件转换为可处理的类别；没有后续县档或物质证据，不能由分类推定实际生活。账本提示的研究线索为宫廷政治、制度运作与中央—地方信息链研究。事件所处条件为城市贫困、治安和宫禁漏洞叠加。后续变化应限于清廷随后整顿城防、旗营与案件处置。来源限制：以嘉庆八年（1803年）的同期文书为定位起点；官修记录、奏报或外交文本服务特定机构立场，具体日期、规模、地方执行与对手视角须以独立材料复核。
\hypertarget{civic:EQ-1804-COASTAL-PIRACY}{}\paragraph{嘉庆九年（1804年）}东南海盗活动加剧，福建、广东沿海军政压力上升。本事件需要放在市场价格、交通路线、水文环境、边贸和地方债务的相互制约中理解，不把中央决策与区域经济割裂开来。 核心取证为《清仁宗实录》嘉庆九年相关条；军机处档、战事方略及地方奏报。编年条目可以为行动建立相对次序，却通常不保存劳动过程、物价变动、建筑材料或非精英的叙述。账本提示的研究线索为战争动员、军需、战场暴力与地方社会研究。事件所处条件为海上贸易、武装网络与官府海防能力失衡。后续变化应限于海盗、商人和地方居民角色不能由官方标签一概而论。来源限制：以嘉庆九年（1804年）的同期文书为定位起点；官修记录、奏报或外交文本服务特定机构立场，具体日期、规模、地方执行与对手视角须以独立材料复核。
\hypertarget{civic:EQ-1805-WHITE-LOTUS-WAR-END}{}\paragraph{嘉庆十年（1805年）}白莲教大规模战争基本结束，清廷转入善后和治安重整。本事件需要放在市场价格、交通路线、水文环境、边贸和地方债务的相互制约中理解，不把中央决策与区域经济割裂开来。 核心取证为《清仁宗实录》嘉庆十年相关条；军机处档、战事方略及地方奏报。编年条目可以为行动建立相对次序，却通常不保存劳动过程、物价变动、建筑材料或非精英的叙述。账本提示的研究线索为战争动员、军需、战场暴力与地方社会研究。事件所处条件为持续清剿、招抚和地方武装动员共同作用。后续变化应限于战后财政、人口与地方武装遗产长期存在。来源限制：以嘉庆十年（1805年）的同期文书为定位起点；官修记录、奏报或外交文本服务特定机构立场，具体日期、规模、地方执行与对手视角须以独立材料复核。
\hypertarget{civic:EQ-1806-TSAI-QIAN-PIRACY}{}\paragraph{嘉庆十一年（1806年）}蔡牵等海盗集团持续活动，清廷调整水师与沿海防务。本事件需要放在市场价格、交通路线、水文环境、边贸和地方债务的相互制约中理解，不把中央决策与区域经济割裂开来。 核心取证为《清仁宗实录》嘉庆十一年相关条；军机处档、战事方略及地方奏报。编年条目可以为行动建立相对次序，却通常不保存劳动过程、物价变动、建筑材料或非精英的叙述。账本提示的研究线索为战争动员、军需、战场暴力与地方社会研究。事件所处条件为沿海贸易网络可为武装集团提供补给与信息。后续变化应限于水师战果和海盗人数须以多方材料审慎核对。来源限制：以嘉庆十一年（1806年）的同期文书为定位起点；官修记录、奏报或外交文本服务特定机构立场，具体日期、规模、地方执行与对手视角须以独立材料复核。
\hypertarget{civic:EQ-1807-COASTAL-NAVAL-REFORM}{}\paragraph{嘉庆十二年（1807年）}清廷继续通过水师调度、招抚和地方协作应对海盗。本事件需要放在市场价格、交通路线、水文环境、边贸和地方债务的相互制约中理解，不把中央决策与区域经济割裂开来。 核心取证为《清仁宗实录》嘉庆十二年相关条；宫中朱批奏折、内阁大库题本与《清史稿》相关志传。行政文书显示的是机构如何把道路、港口、人口、宗教或案件转换为可处理的类别；没有后续县档或物质证据，不能由分类推定实际生活。账本提示的研究线索为宫廷政治、制度运作与中央—地方信息链研究。事件所处条件为传统水师装备与指挥难适应流动海战。后续变化应限于海防改革的成效应以具体港口和海域资料检验。来源限制：以嘉庆十二年（1807年）的同期文书为定位起点；官修记录、奏报或外交文本服务特定机构立场，具体日期、规模、地方执行与对手视角须以独立材料复核。
\hypertarget{civic:EQ-1808-PIRATE-ALLIANCES}{}\paragraph{嘉庆十三年（1808年）}郑一嫂、张保仔等海盗网络在华南海域形成强大联盟。本事件需要放在市场价格、交通路线、水文环境、边贸和地方债务的相互制约中理解，不把中央决策与区域经济割裂开来。 核心取证为《清仁宗实录》嘉庆十三年相关条；地方志、督抚奏折及地方档案。编年条目可以为行动建立相对次序，却通常不保存劳动过程、物价变动、建筑材料或非精英的叙述。账本提示的研究线索为地方档案、迁徙、宗教、族群与社会动员研究。事件所处条件为商业航路、渔盐社会和地方保护关系提供空间。后续变化应限于成员规模和控制范围多为官私叙事的估计。来源限制：以嘉庆十三年（1808年）的同期文书为定位起点；官修记录、奏报或外交文本服务特定机构立场，具体日期、规模、地方执行与对手视角须以独立材料复核。
\hypertarget{civic:EQ-1809-SOUTH-CHINA-SEA-CAMPAIGN}{}\paragraph{嘉庆十四年（1809年）}清军、地方势力与外来船只在华南海域共同影响剿盗行动。本事件需要放在市场价格、交通路线、水文环境、边贸和地方债务的相互制约中理解，不把中央决策与区域经济割裂开来。 核心取证为《清仁宗实录》嘉庆十四年相关条；军机处档、战事方略及地方奏报。编年条目可以为行动建立相对次序，却通常不保存劳动过程、物价变动、建筑材料或非精英的叙述。账本提示的研究线索为战争动员、军需、战场暴力与地方社会研究。事件所处条件为海盗联盟威胁港口、航运与官府威信。后续变化应限于跨国海上互动需结合中英文航海和领事记录。来源限制：以嘉庆十四年（1809年）的同期文书为定位起点；官修记录、奏报或外交文本服务特定机构立场，具体日期、规模、地方执行与对手视角须以独立材料复核。
\subsection{1810—1819：区域网络、制度媒介与环境条件}
这些事件的先后可以由账本定位，但每一项影响都须回到具体区域、材料与行动者，不能被压缩为抽象的国家趋势。
\hypertarget{civic:EQ-1810-ZHANG-BAOZAI-SURRENDER}{}\paragraph{嘉庆十五年（1810年）}张保仔接受招安，华南大海盗联盟瓦解。本事件需要放在市场价格、交通路线、水文环境、边贸和地方债务的相互制约中理解，不把中央决策与区域经济割裂开来。 核心取证为《清仁宗实录》嘉庆十五年相关条；军机处档、战事方略及地方奏报。编年条目可以为行动建立相对次序，却通常不保存劳动过程、物价变动、建筑材料或非精英的叙述。账本提示的研究线索为战争动员、军需、战场暴力与地方社会研究。事件所处条件为清廷以军事实力、谈判和职位安置结合处理海盗。后续变化应限于招安不等于沿海社会问题消失，参与者去向需具体追踪。来源限制：以嘉庆十五年（1810年）的同期文书为定位起点；官修记录、奏报或外交文本服务特定机构立场，具体日期、规模、地方执行与对手视角须以独立材料复核。
\hypertarget{civic:EQ-1811-BANNER-LIVELIHOOD}{}\paragraph{嘉庆十六年（1811年）}八旗生计、债务和人口供养问题持续成为中央与地方的治理议题。本事件需要放在市场价格、交通路线、水文环境、边贸和地方债务的相互制约中理解，不把中央决策与区域经济割裂开来。 核心取证为《清仁宗实录》嘉庆十六年相关条；地方志、督抚奏折及地方档案。编年条目可以为行动建立相对次序，却通常不保存劳动过程、物价变动、建筑材料或非精英的叙述。账本提示的研究线索为地方档案、迁徙、宗教、族群与社会动员研究。事件所处条件为旗饷不足与城市物价、家庭规模变化相互影响。后续变化应限于旗人生活差异很大，不能由政策文书概括整体。来源限制：以嘉庆十六年（1811年）的同期文书为定位起点；官修记录、奏报或外交文本服务特定机构立场，具体日期、规模、地方执行与对手视角须以独立材料复核。
\hypertarget{civic:EQ-1812-TIANLI-PREPARATION}{}\paragraph{嘉庆十七年（1812年）}天理教组织活动在直隶、河南、山东等地被官府逐步察觉。本事件需要放在市场价格、交通路线、水文环境、边贸和地方债务的相互制约中理解，不把中央决策与区域经济割裂开来。 核心取证为《清仁宗实录》嘉庆十七年相关条；地方志、督抚奏折及地方档案。编年条目可以为行动建立相对次序，却通常不保存劳动过程、物价变动、建筑材料或非精英的叙述。账本提示的研究线索为地方档案、迁徙、宗教、族群与社会动员研究。事件所处条件为宗教结社、地方交通和秘密组织网络形成行动条件。后续变化应限于官方案件材料重在追捕，不能完全还原参与者动机。来源限制：以嘉庆十七年（1812年）的同期文书为定位起点；官修记录、奏报或外交文本服务特定机构立场，具体日期、规模、地方执行与对手视角须以独立材料复核。
\hypertarget{civic:EQ-1813-TIANLI-UPRISING}{}\paragraph{嘉庆十八年（1813年）}天理教起事者攻入紫禁城部分区域，京师发生严重安全危机。本事件需要放在市场价格、交通路线、水文环境、边贸和地方债务的相互制约中理解，不把中央决策与区域经济割裂开来。 核心取证为《清仁宗实录》嘉庆十八年相关条；军机处档、战事方略及地方奏报。编年条目可以为行动建立相对次序，却通常不保存劳动过程、物价变动、建筑材料或非精英的叙述。账本提示的研究线索为战争动员、军需、战场暴力与地方社会研究。事件所处条件为基层结社与京城内外联络突破了既有警卫。后续变化应限于清廷加强案件侦缉和地方控制，审讯文本须注意强迫与分类。来源限制：以嘉庆十八年（1813年）的同期文书为定位起点；官修记录、奏报或外交文本服务特定机构立场，具体日期、规模、地方执行与对手视角须以独立材料复核。
\hypertarget{civic:EQ-1814-TIANLI-SUPPRESSION}{}\paragraph{嘉庆十九年（1814年）}天理教起事被镇压，相关人员遭逮捕和处置。本事件需要放在市场价格、交通路线、水文环境、边贸和地方债务的相互制约中理解，不把中央决策与区域经济割裂开来。 核心取证为《清仁宗实录》嘉庆十九年相关条；宫中朱批奏折、内阁大库题本与《清史稿》相关志传。行政文书显示的是机构如何把道路、港口、人口、宗教或案件转换为可处理的类别；没有后续县档或物质证据，不能由分类推定实际生活。账本提示的研究线索为宫廷政治、制度运作与中央—地方信息链研究。事件所处条件为朝廷将京师遇袭视为统治危机。后续变化应限于案件善后扩大了对结社和宗教活动的监控。来源限制：以嘉庆十九年（1814年）的同期文书为定位起点；官修记录、奏报或外交文本服务特定机构立场，具体日期、规模、地方执行与对手视角须以独立材料复核。
\hypertarget{civic:EQ-1815-GRAND-CANAL-REPAIRS}{}\paragraph{嘉庆二十年（1815年）}漕运、河工和仓储问题在战后财政压力下持续受到讨论。本事件需要放在市场价格、交通路线、水文环境、边贸和地方债务的相互制约中理解，不把中央决策与区域经济割裂开来。 核心取证为《清仁宗实录》嘉庆二十年相关条；户部奏销、盐政、河工或海关档案。行政文书显示的是机构如何把道路、港口、人口、宗教或案件转换为可处理的类别；没有后续县档或物质证据，不能由分类推定实际生活。账本提示的研究线索为财政账册、市场网络、灾害环境与区域经济研究。事件所处条件为黄河水势、运河淤塞和漕粮征解相互牵连。后续变化应限于工程文书的计划与实际运输量必须区分。来源限制：以嘉庆二十年（1815年）的同期文书为定位起点；官修记录、奏报或外交文本服务特定机构立场，具体日期、规模、地方执行与对手视角须以独立材料复核。
\hypertarget{civic:EQ-1816-AMHERST-MISSION}{}\paragraph{嘉庆二十一年（1816年）}阿美士德使团来华，礼仪与外交交涉未获突破。本事件需要放在市场价格、交通路线、水文环境、边贸和地方债务的相互制约中理解，不把中央决策与区域经济割裂开来。 核心取证为《清仁宗实录》嘉庆二十一年相关条；《筹办夷务始末》相关年卷与中外照会、条约文本。行政文书显示的是机构如何把道路、港口、人口、宗教或案件转换为可处理的类别；没有后续县档或物质证据，不能由分类推定实际生活。账本提示的研究线索为条约文本、外交档案、口岸社会与国际法实践研究。事件所处条件为英国希望扩大对华贸易与外交接触。后续变化应限于使团叙述反映英方利益，清方礼仪文本也服务朝廷秩序。来源限制：以嘉庆二十一年（1816年）的同期文书为定位起点；官修记录、奏报或外交文本服务特定机构立场，具体日期、规模、地方执行与对手视角须以独立材料复核。
\hypertarget{civic:EQ-1817-SALT-ARREARS}{}\paragraph{嘉庆二十二年（1817年）}盐课、商欠和财政亏空继续困扰中央财政。本事件需要放在市场价格、交通路线、水文环境、边贸和地方债务的相互制约中理解，不把中央决策与区域经济割裂开来。 核心取证为《清仁宗实录》嘉庆二十二年相关条；户部奏销、盐政、河工或海关档案。行政文书显示的是机构如何把道路、港口、人口、宗教或案件转换为可处理的类别；没有后续县档或物质证据，不能由分类推定实际生活。账本提示的研究线索为财政账册、市场网络、灾害环境与区域经济研究。事件所处条件为战费、商业波动与既有盐政制度缺陷并存。后续变化应限于名义课额、实收与欠额不应混为一谈。来源限制：以嘉庆二十二年（1817年）的同期文书为定位起点；官修记录、奏报或外交文本服务特定机构立场，具体日期、规模、地方执行与对手视角须以独立材料复核。
\hypertarget{civic:EQ-1818-FLOOD-RELIEF}{}\paragraph{嘉庆二十三年（1818年）}黄淮水患和赈务成为地方奏报与中央拨款的重要议题。本事件需要放在市场价格、交通路线、水文环境、边贸和地方债务的相互制约中理解，不把中央决策与区域经济割裂开来。 核心取证为《清仁宗实录》嘉庆二十三年相关条；地方志、督抚奏折及地方档案。编年条目可以为行动建立相对次序，却通常不保存劳动过程、物价变动、建筑材料或非精英的叙述。账本提示的研究线索为地方档案、迁徙、宗教、族群与社会动员研究。事件所处条件为河道、堤防、降水与贫困共同塑造灾害。后续变化应限于灾异记录反映申赈机制，不能直接转换为精确死亡数字。来源限制：以嘉庆二十三年（1818年）的同期文书为定位起点；官修记录、奏报或外交文本服务特定机构立场，具体日期、规模、地方执行与对手视角须以独立材料复核。
\hypertarget{civic:EQ-1819-OPIUM-TRADE-GROWTH}{}\paragraph{嘉庆二十四年（1819年）}鸦片贸易增长引起官员对白银流动、海防和社会秩序的担忧。本事件需要放在市场价格、交通路线、水文环境、边贸和地方债务的相互制约中理解，不把中央决策与区域经济割裂开来。 核心取证为《清仁宗实录》嘉庆二十四年相关条；户部奏销、盐政、河工或海关档案。行政文书显示的是机构如何把道路、港口、人口、宗教或案件转换为可处理的类别；没有后续县档或物质证据，不能由分类推定实际生活。账本提示的研究线索为财政账册、市场网络、灾害环境与区域经济研究。事件所处条件为全球贸易、印度鸦片生产与广州贸易制度相互连接。后续变化应限于走私规模需区别官府估计、商人账目和海关统计。来源限制：以嘉庆二十四年（1819年）的同期文书为定位起点；官修记录、奏报或外交文本服务特定机构立场，具体日期、规模、地方执行与对手视角须以独立材料复核。
\subsection{1820—1829：区域网络、制度媒介与环境条件}
这些事件的先后可以由账本定位，但每一项影响都须回到具体区域、材料与行动者，不能被压缩为抽象的国家趋势。
\hypertarget{civic:EQ-1820-DAO-GUANG-ACCESSION}{}\paragraph{嘉庆二十五年（1820年）}道光帝即位，财政、河工、旗务和沿海贸易问题进入新朝。本事件需要放在市场价格、交通路线、水文环境、边贸和地方债务的相互制约中理解，不把中央决策与区域经济割裂开来。 核心取证为《清仁宗实录》嘉庆二十五年相关条；宫中朱批奏折、内阁大库题本与《清史稿》相关志传。行政文书显示的是机构如何把道路、港口、人口、宗教或案件转换为可处理的类别；没有后续县档或物质证据，不能由分类推定实际生活。账本提示的研究线索为宫廷政治、制度运作与中央—地方信息链研究。事件所处条件为嘉庆朝遗留军费和行政压力。后续变化应限于新帝在节用与整顿之间寻求政策平衡。来源限制：以嘉庆二十五年（1820年）的同期文书为定位起点；官修记录、奏报或外交文本服务特定机构立场，具体日期、规模、地方执行与对手视角须以独立材料复核。
\hypertarget{civic:EQ-1821-IMPERIAL-ECONOMY}{}\paragraph{道光元年（1821年）}道光初年继续检讨内务府、官员经费和财政节用。本事件需要放在市场价格、交通路线、水文环境、边贸和地方债务的相互制约中理解，不把中央决策与区域经济割裂开来。 核心取证为《清宣宗实录》道光元年相关条；户部奏销、盐政、河工或海关档案。行政文书显示的是机构如何把道路、港口、人口、宗教或案件转换为可处理的类别；没有后续县档或物质证据，不能由分类推定实际生活。账本提示的研究线索为财政账册、市场网络、灾害环境与区域经济研究。事件所处条件为国库压力与皇室、军政支出并存。后续变化应限于节用谕旨不直接表明各部门实际支出下降。来源限制：以道光元年（1821年）的同期文书为定位起点；官修记录、奏报或外交文本服务特定机构立场，具体日期、规模、地方执行与对手视角须以独立材料复核。
\hypertarget{civic:EQ-1822-YELLOW-RIVER-FLOOD}{}\paragraph{道光二年（1822年）}黄河水患、堤工和灾赈问题反复进入中央议程。本事件需要放在市场价格、交通路线、水文环境、边贸和地方债务的相互制约中理解，不把中央决策与区域经济割裂开来。 核心取证为《清宣宗实录》道光二年相关条；地方志、督抚奏折及地方档案。编年条目可以为行动建立相对次序，却通常不保存劳动过程、物价变动、建筑材料或非精英的叙述。账本提示的研究线索为地方档案、迁徙、宗教、族群与社会动员研究。事件所处条件为河道变迁与沿河人口密集使风险累积。后续变化应限于不同河段和灾民数量必须分开核对。来源限制：以道光二年（1822年）的同期文书为定位起点；官修记录、奏报或外交文本服务特定机构立场，具体日期、规模、地方执行与对手视角须以独立材料复核。
\hypertarget{civic:EQ-1823-BANNER-DEBT}{}\paragraph{道光三年（1823年）}旗人债务、典当和生计问题持续引起官府讨论。本事件需要放在市场价格、交通路线、水文环境、边贸和地方债务的相互制约中理解，不把中央决策与区域经济割裂开来。 核心取证为《清宣宗实录》道光三年相关条；地方志、督抚奏折及地方档案。编年条目可以为行动建立相对次序，却通常不保存劳动过程、物价变动、建筑材料或非精英的叙述。账本提示的研究线索为地方档案、迁徙、宗教、族群与社会动员研究。事件所处条件为固定饷银与城市消费、家庭负担不相匹配。后续变化应限于个案诉讼不能直接代表全部旗人生活。来源限制：以道光三年（1823年）的同期文书为定位起点；官修记录、奏报或外交文本服务特定机构立场，具体日期、规模、地方执行与对手视角须以独立材料复核。
\hypertarget{civic:EQ-1824-GRAND-CANAL-CRISIS}{}\paragraph{道光四年（1824年）}运河水运和漕粮转输面临河工、淤积与财政问题。本事件需要放在市场价格、交通路线、水文环境、边贸和地方债务的相互制约中理解，不把中央决策与区域经济割裂开来。 核心取证为《清宣宗实录》道光四年相关条；户部奏销、盐政、河工或海关档案。行政文书显示的是机构如何把道路、港口、人口、宗教或案件转换为可处理的类别；没有后续县档或物质证据，不能由分类推定实际生活。账本提示的研究线索为财政账册、市场网络、灾害环境与区域经济研究。事件所处条件为漕运体系依赖复杂的河道、仓储和民夫网络。后续变化应限于政策调整与实际粮运路线需以账册和地方志核验。来源限制：以道光四年（1824年）的同期文书为定位起点；官修记录、奏报或外交文本服务特定机构立场，具体日期、规模、地方执行与对手视角须以独立材料复核。
\hypertarget{civic:EQ-1825-SOUTHWEST-MINING-UNREST}{}\paragraph{道光五年（1825年）}云贵矿业、移民和治安纠纷继续考验地方治理。本事件需要放在市场价格、交通路线、水文环境、边贸和地方债务的相互制约中理解，不把中央决策与区域经济割裂开来。 核心取证为《清宣宗实录》道光五年相关条；地方志、督抚奏折及地方档案。编年条目可以为行动建立相对次序，却通常不保存劳动过程、物价变动、建筑材料或非精英的叙述。账本提示的研究线索为地方档案、迁徙、宗教、族群与社会动员研究。事件所处条件为资源开发吸引流动人口并放大既有权力冲突。后续变化应限于“矿匪”等官方称呼不能覆盖参与者差异。来源限制：以道光五年（1825年）的同期文书为定位起点；官修记录、奏报或外交文本服务特定机构立场，具体日期、规模、地方执行与对手视角须以独立材料复核。
\hypertarget{civic:EQ-1826-JAHANGIR-KHOJA-UPRISING}{}\paragraph{道光六年（1826年）}张格尔和卓进入喀什噶尔，天山南路爆发大规模反清动荡。本事件需要放在市场价格、交通路线、水文环境、边贸和地方债务的相互制约中理解，不把中央决策与区域经济割裂开来。 核心取证为《清宣宗实录》道光六年相关条；军机处档、战事方略及地方奏报。编年条目可以为行动建立相对次序，却通常不保存劳动过程、物价变动、建筑材料或非精英的叙述。账本提示的研究线索为战争动员、军需、战场暴力与地方社会研究。事件所处条件为浩罕、绿洲政治与清军治理矛盾相互作用。后续变化应限于清方“叛乱”叙述须同中亚和地方材料互证。来源限制：以道光六年（1826年）的同期文书为定位起点；官修记录、奏报或外交文本服务特定机构立场，具体日期、规模、地方执行与对手视角须以独立材料复核。
\hypertarget{civic:EQ-1827-KASHGAR-RECOVERY}{}\paragraph{道光七年（1827年）}清军反攻并收复喀什噶尔等地，张格尔势力受挫。本事件需要放在市场价格、交通路线、水文环境、边贸和地方债务的相互制约中理解，不把中央决策与区域经济割裂开来。 核心取证为《清宣宗实录》道光七年相关条；军机处档、战事方略及地方奏报。编年条目可以为行动建立相对次序，却通常不保存劳动过程、物价变动、建筑材料或非精英的叙述。账本提示的研究线索为战争动员、军需、战场暴力与地方社会研究。事件所处条件为清廷集中西北兵力并依赖地方补给。后续变化应限于恢复城镇不等于绿洲治理和贸易网络立即稳定。来源限制：以道光七年（1827年）的同期文书为定位起点；官修记录、奏报或外交文本服务特定机构立场，具体日期、规模、地方执行与对手视角须以独立材料复核。
\hypertarget{civic:EQ-1828-JAHANGIR-CAPTURE}{}\paragraph{道光八年（1828年）}张格尔被捕并押解处置，清廷加强南疆军政控制。本事件需要放在市场价格、交通路线、水文环境、边贸和地方债务的相互制约中理解，不把中央决策与区域经济割裂开来。 核心取证为《清宣宗实录》道光八年相关条；军机处档、理藩院档或相关方略。编年条目可以为行动建立相对次序，却通常不保存劳动过程、物价变动、建筑材料或非精英的叙述。账本提示的研究线索为多语边疆档案、盟旗/寺院/绿洲地方史与帝国治理研究。事件所处条件为军事反攻压缩其活动空间。后续变化应限于案件、处决与地方反应需区分不同材料层次。来源限制：以道光八年（1828年）的同期文书为定位起点；官修记录、奏报或外交文本服务特定机构立场，具体日期、规模、地方执行与对手视角须以独立材料复核。
\hypertarget{civic:EQ-1829-XINJIANG-FINANCE}{}\paragraph{道光九年（1829年）}南疆战后军费、驻防与贸易秩序成为清廷边务财政负担。本事件需要放在市场价格、交通路线、水文环境、边贸和地方债务的相互制约中理解，不把中央决策与区域经济割裂开来。 核心取证为《清宣宗实录》道光九年相关条；户部奏销、盐政、河工或海关档案。行政文书显示的是机构如何把道路、港口、人口、宗教或案件转换为可处理的类别；没有后续县档或物质证据，不能由分类推定实际生活。账本提示的研究线索为财政账册、市场网络、灾害环境与区域经济研究。事件所处条件为远距离驻军和绿洲行政需要持续资源。后续变化应限于预算、实收和地方供给不可由单一军费数额概括。来源限制：以道光九年（1829年）的同期文书为定位起点；官修记录、奏报或外交文本服务特定机构立场，具体日期、规模、地方执行与对手视角须以独立材料复核。
\subsection{1830—1839：区域网络、制度媒介与环境条件}
这些事件的先后可以由账本定位，但每一项影响都须回到具体区域、材料与行动者，不能被压缩为抽象的国家趋势。
\hypertarget{civic:EQ-1830-RUSSO-QING-TRADE}{}\paragraph{道光十年（1830年）}恰克图等边境贸易和俄清交涉持续影响北部边务。本事件需要放在市场价格、交通路线、水文环境、边贸和地方债务的相互制约中理解，不把中央决策与区域经济割裂开来。 核心取证为《清宣宗实录》道光十年相关条；《筹办夷务始末》相关年卷与中外照会、条约文本。行政文书显示的是机构如何把道路、港口、人口、宗教或案件转换为可处理的类别；没有后续县档或物质证据，不能由分类推定实际生活。账本提示的研究线索为条约文本、外交档案、口岸社会与国际法实践研究。事件所处条件为跨境商业与边界管理相互依存。后续变化应限于贸易额和走私规模需要双方档案互校。来源限制：以道光十年（1830年）的同期文书为定位起点；官修记录、奏报或外交文本服务特定机构立场，具体日期、规模、地方执行与对手视角须以独立材料复核。
\hypertarget{civic:EQ-1831-CENTRAL-ASIA-CHOLERA}{}\paragraph{道光十一年（1831年）}疾病流行与交通网络对军队、城市和商路造成压力。本事件需要放在市场价格、交通路线、水文环境、边贸和地方债务的相互制约中理解，不把中央决策与区域经济割裂开来。 核心取证为《清宣宗实录》道光十一年相关条；地方志、督抚奏折及地方档案。编年条目可以为行动建立相对次序，却通常不保存劳动过程、物价变动、建筑材料或非精英的叙述。账本提示的研究线索为地方档案、迁徙、宗教、族群与社会动员研究。事件所处条件为区域流动和环境条件影响疾病传播。后续变化应限于病例和死亡数据受记录制度限制，不能过度精确。来源限制：以道光十一年（1831年）的同期文书为定位起点；官修记录、奏报或外交文本服务特定机构立场，具体日期、规模、地方执行与对手视角须以独立材料复核。
\hypertarget{civic:EQ-1832-COASTAL-DEFENCE}{}\paragraph{道光十二年（1832年）}鸦片走私增长背景下的海防、缉私和地方官责任成为争议。本事件需要放在市场价格、交通路线、水文环境、边贸和地方债务的相互制约中理解，不把中央决策与区域经济割裂开来。 核心取证为《清宣宗实录》道光十二年相关条；宫中朱批奏折、内阁大库题本与《清史稿》相关志传。行政文书显示的是机构如何把道路、港口、人口、宗教或案件转换为可处理的类别；没有后续县档或物质证据，不能由分类推定实际生活。账本提示的研究线索为宫廷政治、制度运作与中央—地方信息链研究。事件所处条件为广州贸易制度难以控制海上走私网络。后续变化应限于禁令与缉私战报不能直接证明贸易被阻断。来源限制：以道光十二年（1832年）的同期文书为定位起点；官修记录、奏报或外交文本服务特定机构立场，具体日期、规模、地方执行与对手视角须以独立材料复核。
\hypertarget{civic:EQ-1833-OPIUM-POLICY-DEBATE}{}\paragraph{道光十三年（1833年）}官员围绕禁烟、征税、白银和海防提出相互竞争的方案。本事件需要放在市场价格、交通路线、水文环境、边贸和地方债务的相互制约中理解，不把中央决策与区域经济割裂开来。 核心取证为《清宣宗实录》道光十三年相关条；宫中朱批奏折、内阁大库题本与《清史稿》相关志传。行政文书显示的是机构如何把道路、港口、人口、宗教或案件转换为可处理的类别；没有后续县档或物质证据，不能由分类推定实际生活。账本提示的研究线索为宫廷政治、制度运作与中央—地方信息链研究。事件所处条件为鸦片输入、银价波动和行政腐败引发担忧。后续变化应限于奏议反映立场，不等于朝廷形成一致政策。来源限制：以道光十三年（1833年）的同期文书为定位起点；官修记录、奏报或外交文本服务特定机构立场，具体日期、规模、地方执行与对手视角须以独立材料复核。
\hypertarget{civic:EQ-1834-NAPIER-AFFAIR}{}\paragraph{道光十四年（1834年）}律劳卑抵广州后与两广官员发生外交摩擦。本事件需要放在市场价格、交通路线、水文环境、边贸和地方债务的相互制约中理解，不把中央决策与区域经济割裂开来。 核心取证为《清宣宗实录》道光十四年相关条；《筹办夷务始末》相关年卷与中外照会、条约文本。行政文书显示的是机构如何把道路、港口、人口、宗教或案件转换为可处理的类别；没有后续县档或物质证据，不能由分类推定实际生活。账本提示的研究线索为条约文本、外交档案、口岸社会与国际法实践研究。事件所处条件为英国希望绕开广州商馆制度直接交涉。后续变化应限于双方文书的称谓和程序冲突不应简化为文化误解。来源限制：以道光十四年（1834年）的同期文书为定位起点；官修记录、奏报或外交文本服务特定机构立场，具体日期、规模、地方执行与对手视角须以独立材料复核。
\hypertarget{civic:EQ-1835-GUANGZHOU-TRADE-REGULATION}{}\paragraph{道光十五年（1835年）}广州贸易、行商债务和外商管理继续在地方与中央之间协商。本事件需要放在市场价格、交通路线、水文环境、边贸和地方债务的相互制约中理解，不把中央决策与区域经济割裂开来。 核心取证为《清宣宗实录》道光十五年相关条；户部奏销、盐政、河工或海关档案。行政文书显示的是机构如何把道路、港口、人口、宗教或案件转换为可处理的类别；没有后续县档或物质证据，不能由分类推定实际生活。账本提示的研究线索为财政账册、市场网络、灾害环境与区域经济研究。事件所处条件为一口贸易下的税收、信用和走私问题累积。后续变化应限于行商档案偏向商人和官府，需与外文商业记录互读。来源限制：以道光十五年（1835年）的同期文书为定位起点；官修记录、奏报或外交文本服务特定机构立场，具体日期、规模、地方执行与对手视角须以独立材料复核。
\hypertarget{civic:EQ-1836-OPIUM-LEGALIZATION-DEBATE}{}\paragraph{道光十六年（1836年）}许乃济等围绕鸦片弛禁、征税与禁绝提出政策论争。本事件需要放在市场价格、交通路线、水文环境、边贸和地方债务的相互制约中理解，不把中央决策与区域经济割裂开来。 核心取证为《清宣宗实录》道光十六年相关条；宫中朱批奏折、内阁大库题本与《清史稿》相关志传。行政文书显示的是机构如何把道路、港口、人口、宗教或案件转换为可处理的类别；没有后续县档或物质证据，不能由分类推定实际生活。账本提示的研究线索为宫廷政治、制度运作与中央—地方信息链研究。事件所处条件为禁令难以遏制走私且银价问题突出。后续变化应限于奏议是方案而非实施结果，须同随后禁烟政策区分。来源限制：以道光十六年（1836年）的同期文书为定位起点；官修记录、奏报或外交文本服务特定机构立场，具体日期、规模、地方执行与对手视角须以独立材料复核。
\hypertarget{civic:EQ-1837-CANTON-SMUGGLING}{}\paragraph{道光十七年（1837年）}广东缉私与鸦片走私之间的冲突持续升级。本事件需要放在市场价格、交通路线、水文环境、边贸和地方债务的相互制约中理解，不把中央决策与区域经济割裂开来。 核心取证为《清宣宗实录》道光十七年相关条；地方志、督抚奏折及地方档案。编年条目可以为行动建立相对次序，却通常不保存劳动过程、物价变动、建筑材料或非精英的叙述。账本提示的研究线索为地方档案、迁徙、宗教、族群与社会动员研究。事件所处条件为沿海商人、船户、官员与外商形成复杂利益网络。后续变化应限于查获数字不等于走私总量。来源限制：以道光十七年（1837年）的同期文书为定位起点；官修记录、奏报或外交文本服务特定机构立场，具体日期、规模、地方执行与对手视角须以独立材料复核。
\hypertarget{civic:EQ-1838-LIN-ZEXU-APPOINTMENT}{}\paragraph{道光十八年（1838年）}林则徐受命赴广东查禁鸦片，禁烟政策转为高压执行。本事件需要放在市场价格、交通路线、水文环境、边贸和地方债务的相互制约中理解，不把中央决策与区域经济割裂开来。 核心取证为《清宣宗实录》道光十八年相关条；宫中朱批奏折、内阁大库题本与《清史稿》相关志传。行政文书显示的是机构如何把道路、港口、人口、宗教或案件转换为可处理的类别；没有后续县档或物质证据，不能由分类推定实际生活。账本提示的研究线索为宫廷政治、制度运作与中央—地方信息链研究。事件所处条件为道光帝决定以严禁回应贸易和财政危机。后续变化应限于任命文书表明政策方向，地方执行与外商反应另需追踪。来源限制：以道光十八年（1838年）的同期文书为定位起点；官修记录、奏报或外交文本服务特定机构立场，具体日期、规模、地方执行与对手视角须以独立材料复核。
\hypertarget{civic:EQ-1839-HUMEN-OPIUM-DESTRUCTION}{}\paragraph{道光十九年（1839年）}林则徐在虎门销毁收缴鸦片，禁烟冲突公开化。本事件需要放在市场价格、交通路线、水文环境、边贸和地方债务的相互制约中理解，不把中央决策与区域经济割裂开来。 核心取证为《清宣宗实录》道光十九年相关条；《筹办夷务始末》相关年卷与中外照会、条约文本。行政文书显示的是机构如何把道路、港口、人口、宗教或案件转换为可处理的类别；没有后续县档或物质证据，不能由分类推定实际生活。账本提示的研究线索为条约文本、外交档案、口岸社会与国际法实践研究。事件所处条件为广东禁烟迫使外商交出库存并改变贸易秩序。后续变化应限于数量、所有权和各方责任需核对中英档案。来源限制：以道光十九年（1839年）的同期文书为定位起点；官修记录、奏报或外交文本服务特定机构立场，具体日期、规模、地方执行与对手视角须以独立材料复核。
\subsection{1840—1849：区域网络、制度媒介与环境条件}
这些事件的先后可以由账本定位，但每一项影响都须回到具体区域、材料与行动者，不能被压缩为抽象的国家趋势。
\hypertarget{civic:EQ-1840-FIRST-OPIUM-WAR}{}\paragraph{道光二十年（1840年）}英国舰队来华，第一次鸦片战争爆发。本事件需要放在市场价格、交通路线、水文环境、边贸和地方债务的相互制约中理解，不把中央决策与区域经济割裂开来。 核心取证为《清宣宗实录》道光二十年相关条；军机处档、战事方略及地方奏报。编年条目可以为行动建立相对次序，却通常不保存劳动过程、物价变动、建筑材料或非精英的叙述。账本提示的研究线索为战争动员、军需、战场暴力与地方社会研究。事件所处条件为禁烟冲突与英国以武力重塑贸易外交关系相结合。后续变化应限于战事涉及多口岸与地方社会，不能只以条约谈判叙述。来源限制：以道光二十年（1840年）的同期文书为定位起点；官修记录、奏报或外交文本服务特定机构立场，具体日期、规模、地方执行与对手视角须以独立材料复核。
\hypertarget{civic:EQ-1841-CHUANBI-NEGOTIATIONS}{}\paragraph{道光二十一年（1841年）}穿鼻草约等谈判和沿海战事并行，清廷内部对和战意见分歧。本事件需要放在市场价格、交通路线、水文环境、边贸和地方债务的相互制约中理解，不把中央决策与区域经济割裂开来。 核心取证为《清宣宗实录》道光二十一年相关条；《筹办夷务始末》相关年卷与中外照会、条约文本。行政文书显示的是机构如何把道路、港口、人口、宗教或案件转换为可处理的类别；没有后续县档或物质证据，不能由分类推定实际生活。账本提示的研究线索为条约文本、外交档案、口岸社会与国际法实践研究。事件所处条件为军事压力迫使地方官与英方接触。后续变化应限于草约地位、授权和后续否决须分别说明。来源限制：以道光二十一年（1841年）的同期文书为定位起点；官修记录、奏报或外交文本服务特定机构立场，具体日期、规模、地方执行与对手视角须以独立材料复核。
\hypertarget{civic:EQ-1841-CANTON-BATTLES}{}\paragraph{道光二十一年（1841年）}广州周边战事和赔款谈判反复进行。本事件需要放在市场价格、交通路线、水文环境、边贸和地方债务的相互制约中理解，不把中央决策与区域经济割裂开来。 核心取证为《清宣宗实录》道光二十一年相关条；军机处档、战事方略及地方奏报。编年条目可以为行动建立相对次序，却通常不保存劳动过程、物价变动、建筑材料或非精英的叙述。账本提示的研究线索为战争动员、军需、战场暴力与地方社会研究。事件所处条件为珠江口炮台、商贸城市与英军海军优势相互作用。后续变化应限于战果、平民损失与地方抗战需多源核验。来源限制：以道光二十一年（1841年）的同期文书为定位起点；官修记录、奏报或外交文本服务特定机构立场，具体日期、规模、地方执行与对手视角须以独立材料复核。
\hypertarget{civic:EQ-1842-NANJING-TREATY}{}\paragraph{道光二十二年（1842年）}《南京条约》签订，清廷割让香港并开放五口通商、支付赔款。本事件需要放在市场价格、交通路线、水文环境、边贸和地方债务的相互制约中理解，不把中央决策与区域经济割裂开来。 核心取证为《清宣宗实录》道光二十二年相关条；《筹办夷务始末》相关年卷与中外照会、条约文本。行政文书显示的是机构如何把道路、港口、人口、宗教或案件转换为可处理的类别；没有后续县档或物质证据，不能由分类推定实际生活。账本提示的研究线索为条约文本、外交档案、口岸社会与国际法实践研究。事件所处条件为清军失利和长江航运受威胁迫使议和。后续变化应限于签署、批准、换约和各口岸执行是不同阶段。来源限制：以道光二十二年（1842年）的同期文书为定位起点；官修记录、奏报或外交文本服务特定机构立场，具体日期、规模、地方执行与对手视角须以独立材料复核。
\hypertarget{civic:EQ-1843-BOGUE-TREATY}{}\paragraph{道光二十三年（1843年）}《虎门条约》补充领事裁判、协定关税等安排。本事件需要放在市场价格、交通路线、水文环境、边贸和地方债务的相互制约中理解，不把中央决策与区域经济割裂开来。 核心取证为《清宣宗实录》道光二十三年相关条；《筹办夷务始末》相关年卷与中外照会、条约文本。行政文书显示的是机构如何把道路、港口、人口、宗教或案件转换为可处理的类别；没有后续县档或物质证据，不能由分类推定实际生活。账本提示的研究线索为条约文本、外交档案、口岸社会与国际法实践研究。事件所处条件为南京条约留下具体通商制度未定问题。后续变化应限于中英文文本及其适用范围需并读。来源限制：以道光二十三年（1843年）的同期文书为定位起点；官修记录、奏报或外交文本服务特定机构立场，具体日期、规模、地方执行与对手视角须以独立材料复核。
\hypertarget{civic:EQ-1844-WANGXIA-WHAMPOA-TREATIES}{}\paragraph{道光二十四年（1844年）}《望厦条约》《黄埔条约》分别确立美法在华通商与相关权利。本事件需要放在市场价格、交通路线、水文环境、边贸和地方债务的相互制约中理解，不把中央决策与区域经济割裂开来。 核心取证为《清宣宗实录》道光二十四年相关条；《筹办夷务始末》相关年卷与中外照会、条约文本。行政文书显示的是机构如何把道路、港口、人口、宗教或案件转换为可处理的类别；没有后续县档或物质证据，不能由分类推定实际生活。账本提示的研究线索为条约文本、外交档案、口岸社会与国际法实践研究。事件所处条件为英清条约改变口岸制度并吸引其他列强跟进。后续变化应限于最惠国、传教与领事权的执行须以地方案件验证。来源限制：以道光二十四年（1844年）的同期文书为定位起点；官修记录、奏报或外交文本服务特定机构立场，具体日期、规模、地方执行与对手视角须以独立材料复核。
\hypertarget{civic:EQ-1845-TREATY-PORT-OPENING}{}\paragraph{道光二十五年（1845年）}五口通商制度进入具体开埠、税则和城市管理阶段。本事件需要放在市场价格、交通路线、水文环境、边贸和地方债务的相互制约中理解，不把中央决策与区域经济割裂开来。 核心取证为《清宣宗实录》道光二十五年相关条；户部奏销、盐政、河工或海关档案。行政文书显示的是机构如何把道路、港口、人口、宗教或案件转换为可处理的类别；没有后续县档或物质证据，不能由分类推定实际生活。账本提示的研究线索为财政账册、市场网络、灾害环境与区域经济研究。事件所处条件为条约文本需要转化为海关、领事和地方行政程序。后续变化应限于开港不等于城市空间和贸易立即按文本运作。来源限制：以道光二十五年（1845年）的同期文书为定位起点；官修记录、奏报或外交文本服务特定机构立场，具体日期、规模、地方执行与对手视角须以独立材料复核。
\hypertarget{civic:EQ-1846-MISSIONARY-RETURN}{}\paragraph{道光二十六年（1846年）}条约环境下传教活动与地方教案问题逐渐增多。本事件需要放在市场价格、交通路线、水文环境、边贸和地方债务的相互制约中理解，不把中央决策与区域经济割裂开来。 核心取证为《清宣宗实录》道光二十六年相关条；地方志、督抚奏折及地方档案。编年条目可以为行动建立相对次序，却通常不保存劳动过程、物价变动、建筑材料或非精英的叙述。账本提示的研究线索为地方档案、迁徙、宗教、族群与社会动员研究。事件所处条件为宗教禁令、条约解释和地方土地纠纷相互交织。后续变化应限于传教士、教民与地方官的叙事都有明确立场。来源限制：以道光二十六年（1846年）的同期文书为定位起点；官修记录、奏报或外交文本服务特定机构立场，具体日期、规模、地方执行与对手视角须以独立材料复核。
\hypertarget{civic:EQ-1847-GUANGZHOU-CONFLICT}{}\paragraph{道光二十七年（1847年）}广州周边的中英冲突显示条约执行和城市主权问题尚未解决。本事件需要放在市场价格、交通路线、水文环境、边贸和地方债务的相互制约中理解，不把中央决策与区域经济割裂开来。 核心取证为《清宣宗实录》道光二十七年相关条；《筹办夷务始末》相关年卷与中外照会、条约文本。行政文书显示的是机构如何把道路、港口、人口、宗教或案件转换为可处理的类别；没有后续县档或物质证据，不能由分类推定实际生活。账本提示的研究线索为条约文本、外交档案、口岸社会与国际法实践研究。事件所处条件为英国要求扩大进入广州的权利与地方抵触相冲突。后续变化应限于地方冲突与中央外交方案需分别分析。来源限制：以道光二十七年（1847年）的同期文书为定位起点；官修记录、奏报或外交文本服务特定机构立场，具体日期、规模、地方执行与对手视角须以独立材料复核。
\hypertarget{civic:EQ-1848-YANGTZE-TRADE}{}\paragraph{道光二十八年（1848年）}通商口岸贸易网络向长江和沿海延伸，地方商人开始适应新制度。本事件需要放在市场价格、交通路线、水文环境、边贸和地方债务的相互制约中理解，不把中央决策与区域经济割裂开来。 核心取证为《清宣宗实录》道光二十八年相关条；户部奏销、盐政、河工或海关档案。行政文书显示的是机构如何把道路、港口、人口、宗教或案件转换为可处理的类别；没有后续县档或物质证据，不能由分类推定实际生活。账本提示的研究线索为财政账册、市场网络、灾害环境与区域经济研究。事件所处条件为关税、航运和外商资本改变既有商业联系。后续变化应限于海关及外贸统计仅覆盖制度内流通。来源限制：以道光二十八年（1848年）的同期文书为定位起点；官修记录、奏报或外交文本服务特定机构立场，具体日期、规模、地方执行与对手视角须以独立材料复核。
\hypertarget{civic:EQ-1849-GUANGZHOU-CITY-ENTRY}{}\paragraph{道光二十九年（1849年）}英国要求进入广州城的争议再次激化，地方民众与官府关系受影响。本事件需要放在市场价格、交通路线、水文环境、边贸和地方债务的相互制约中理解，不把中央决策与区域经济割裂开来。 核心取证为《清宣宗实录》道光二十九年相关条；地方志、督抚奏折及地方档案。编年条目可以为行动建立相对次序，却通常不保存劳动过程、物价变动、建筑材料或非精英的叙述。账本提示的研究线索为地方档案、迁徙、宗教、族群与社会动员研究。事件所处条件为条约口岸权利的解释与城市空间控制冲突。后续变化应限于群众动员的规模和自发性需审慎核验。来源限制：以道光二十九年（1849年）的同期文书为定位起点；官修记录、奏报或外交文本服务特定机构立场，具体日期、规模、地方执行与对手视角须以独立材料复核。
\subsection{1850—1859：区域网络、制度媒介与环境条件}
这些事件的先后可以由账本定位，但每一项影响都须回到具体区域、材料与行动者，不能被压缩为抽象的国家趋势。
\hypertarget{civic:EQ-1850-XIANFENG-ACCESSION}{}\paragraph{道光三十年（1850年）}咸丰帝即位，财政紧张、内乱和对外压力同时加重。本事件需要放在市场价格、交通路线、水文环境、边贸和地方债务的相互制约中理解，不把中央决策与区域经济割裂开来。 核心取证为《清宣宗实录》道光三十年相关条；宫中朱批奏折、内阁大库题本与《清史稿》相关志传。行政文书显示的是机构如何把道路、港口、人口、宗教或案件转换为可处理的类别；没有后续县档或物质证据，不能由分类推定实际生活。账本提示的研究线索为宫廷政治、制度运作与中央—地方信息链研究。事件所处条件为道光帝去世时内外危机已累积。后续变化应限于新朝很快面对太平天国和第二次鸦片战争。来源限制：以道光三十年（1850年）的同期文书为定位起点；官修记录、奏报或外交文本服务特定机构立场，具体日期、规模、地方执行与对手视角须以独立材料复核。
